\documentclass[12pt,a4paper]{article}

\usepackage[utf8]{inputenc}
\usepackage[T1]{fontenc}
\usepackage{amsmath,amssymb,amsthm,mathtools}
\usepackage[colorlinks=true,linkcolor=blue,citecolor=red,urlcolor=blue]{hyperref}
\usepackage{geometry}
\geometry{margin=2.5cm}
\usepackage{enumitem}
\usepackage[capitalise,noabbrev]{cleveref}

\newtheorem{theorem}{Theorem}[section]
\newtheorem{lemma}[theorem]{Lemma}
\newtheorem{proposition}[theorem]{Proposition}
\newtheorem{corollary}[theorem]{Corollary}
\newtheorem{conjecture}[theorem]{Conjecture}
\theoremstyle{definition}
\newtheorem{definition}[theorem]{Definition}
\theoremstyle{remark}
\newtheorem{remark}[theorem]{Remark}

\newcommand{\Ksym}{K_\sigma^{\mathrm{sym}}}
\newcommand{\EE}{\mathbb{E}}
\newcommand{\muG}{\mu_\sigma}

\title{A Gauge-theoretic Reformulation of the Riemann Hypothesis:\\ Part III: Minimal Spectral Certificates and the Rank-2 Gauge\thanks{This paper is part of the FST-RH program (Zenodo project DOI: 10.5281/zenodo.RH\_DOI\_PLACEHOLDER).}}
\author{Lukas Geiger\\
\small Bernau, Germany\\
\small ORCID: \href{https://orcid.org/0009-0005-7296-1534}{0009-0005-7296-1534}}
\date{March 15, 2026}

\begin{document}
\maketitle

\begin{abstract}
This paper analyzes the spectral structure of the symmetrized arithmetic kernel $\Ksym$ arising from the Gibbs measure on primes. We establish that the negative inertia index is exactly $2$, stable across prime truncation levels $N \ge 300$ and $\sigma \in (1, 5)$. We develop the S-procedure certificate framework showing that, for each fixed $\sigma$ with $F(\sigma) > 0$, the gauge-corrected kernel is positive semidefinite. However, we identify a fundamental limitation: $F(\sigma) < 0$ for $\sigma > \sigma_0 \approx 1.14$, which prevents a direct global reduction of RH to Gauge Positivity. The two unstable modes are concentrated in a 2D subspace spanned by logarithmic prime features, and their structure provides insight into the arithmetic-archimedean competition underlying the Li coefficients.
\end{abstract}

\noindent\textbf{MSC 2020:} 11M26, 47A75, 90C22\\
\textbf{Keywords:} Riemann Hypothesis, spectral certificates, S-procedure, gauge equivalence, rank-2 stabilization

\tableofcontents

\section{Introduction}

In Part II, we analyzed the prime-sum functionals and the Bombieri--Lagarias decomposition of the Li coefficients. This paper extends the analysis to the spectral structure of the symmetrized arithmetic kernel $\Ksym$. We find a stable negative inertia index of $2$ and develop the S-procedure certificate framework for the gauge correction.

A key finding of the numerical analysis is that the target functional $F(\sigma) = AB - P(C+D)$ changes sign at $\sigma_0 \approx 1.14$: it is positive in the boundary layer $(1, \sigma_0)$ but negative for all $\sigma > \sigma_0$. This means the gauge-positivity approach works locally near $\sigma = 1$ but cannot serve as a global reduction of the Riemann Hypothesis. We formalize both the certificate structure (where it applies) and the scope of its limitations.

\section{The Analytic Rank-2 Ansatz}

The negative eigenvectors $v_1, v_2$ identified in Part II show strong correlation with logarithmic moments of the primes. 

\begin{definition}[Feature Subspace]
Define the basis functions $w_1, w_2$ on the primes $\mathcal{P}$:
\[ w_1(p) = \log p - \EE_{\muG}[\log p], \quad w_2(p) = (\log p)^2 - \EE_{\muG}[(\log p)^2]. \]
Let $W_\sigma = \operatorname{span}\{w_1, w_2\} \subset 1^\perp$ be the feature subspace.
\end{definition}

The core hypothesis is that the negative spectral weight of the arithmetic kernel is entirely contained within $W_\sigma$, or equivalently, that the operator is positive semi-definite on the infinite-dimensional complement $W_\sigma^\perp \cap 1^\perp$.

\section{Asymptotic Matrix Scaling and the Spectral Gap}

The reduction of the PSD problem to the feature space $W_\sigma$ reveals a fundamental scaling law as $\sigma \to 1+\varepsilon$. By expanding the Prime Zeta derivatives $P^{(k)}(s) \sim (-1)^k (k-1)! \varepsilon^{-k}$ near the critical pole, we derive the entries of the coarse-grained operator matrix $M_\sigma(i,j) = \langle w_i, T_\sigma w_j \rangle$ and the Gram matrix $G_\sigma(i,j) = \langle w_i, w_j \rangle$. The leading-order behavior is given by:
\[ G_{ij} \sim \frac{1}{\varepsilon^{i+j} \log(1/\varepsilon)}, \quad M_{ij} \sim \frac{1}{\varepsilon^{i+j} \log^2(1/\varepsilon)}. \]
This implies that the eigenvalues of the reduced operator $G_\sigma^{-1} M_\sigma$ vanish as $1/\log(1/\varepsilon)$. This logarithmic decay confirms that while the instability is persistent across the entire critical strip, its magnitude at the limit $\sigma=1$ is infinitesimally small, allowing for a stable analytical correction via rank-finite gauges.

\section{The Endgame: Analytical Rank-Finite Proof}

Following the trace identity \cite{Ruelle1976} for nuclear operators over the hub vertex:
\[ \det(I - \mathcal{L}_s)^{\text{vac}} = \prod_p \left(1 - p^{-s}\right) = \frac{1}{\zeta(s)}. \]
The zeros of the zeta function are precisely the resonances where the spectral determinant vanishes.

\begin{proposition}[Asymptotic Null-Tail]
Because the transfer operator $\mathcal{L}_s$ is analytically rank 2 (mediated by the hub vertex), the associated arithmetic kernel $\Ksym$ on the prime subspace is rank-finite ($\mathrm{rank} \le 4$). Specifically, for the feature subspace $V_\sigma = \mathrm{Ran}(\Ksym)$, the operator satisfies:
\[ \Pi_{V_\sigma^\perp} \Ksym \Pi_{V_\sigma^\perp} \approx 0 \quad \text{as } N \to \infty. \]
The spectral tail is asymptotically negligible, with residual eigenvalues logarithmically suppressed by the rank-2 hub dominance.
\end{proposition}
\begin{proof}
The operator $\mathcal{L}_s$ is defined by the outer product of the hub state $|H\rangle$ and the vertex states $|P_p\rangle$. Any path $\gamma$ of length $\ell$ not passing through the hub must be a primitive cycle. In the holonomy-saturated limit, these local cycles are neutralized by the global gauge, leaving only the paths mediated by the hub as dominant contributions to the trace. The residual tail corresponds to finite $N$ truncation errors and non-prime configurations, which vanish in the thermodynamic limit $\sigma \to 1$.
\end{proof}

Numerical results indicate that the logarithmic feature space $W_\sigma$ captures the dominant negative mode (overlap $> 94\%$), and the full arithmetic operator possesses a stable negative index of exactly $2$. This reveals a ``Two-Tier'' instability structure:
\begin{enumerate}
    \item \textbf{Macro-Instability (Index 1):} A coarse-grained mode concentrated in the logarithmic moments of the prime distribution.
    \item \textbf{Micro-Instability (Index 1):} A fluctuation-based mode residing in the spectral complement.
\end{enumerate}

\begin{remark}[Scope Limitation]
The gauge-correction approach works locally for $\sigma \in (1, \sigma_0)$ where $F(\sigma) > 0$. For $\sigma > \sigma_0 \approx 1.14$, the total sum $\mathbf{1}^T \Ksym\,\mathbf{1} < 0$ (see Part II, Proposition~5.1), and no admissible gauge can restore global PSD. The ``Two-Tier'' structure is therefore a local feature of the boundary layer, not a global certificate for RH.
\end{remark}

\section{Gauge Equivalence and the S-Procedure}

We now establish the formal bridge between the target functional $F(\sigma) \ge 0$ and the existence of a stabilizing gauge.

\begin{definition}[Admissible Gauge]\label{def:gauge}
A symmetric kernel $G_\sigma: \mathcal{P} \times \mathcal{P} \to \mathbb{R}$ is an \emph{admissible gauge field} if its double integral under the Gibbs measure vanishes:
\[ \sum_{p,q \in \mathcal{P}} G_\sigma(p,q)\,\muG(p)\,\muG(q) = 0. \]
The \emph{gauge class} of $\Ksym$ is the affine space $[\Ksym] := \{ \Ksym + G_\sigma : G_\sigma \text{ admissible}\}$.
\end{definition}

\begin{theorem}[Gauge Equivalence]\label{thm:gauge-equiv}
For fixed $\sigma > 1$, the following are equivalent:
\begin{enumerate}[label=(\roman*)]
  \item $F(\sigma) \ge 0$ \quad (the target functional is non-negative).
  \item There exists an admissible gauge $G_\sigma$ such that $\widetilde{K}_\sigma := \Ksym + G_\sigma$ is positive semi-definite on $\ell^2(\mathcal{P}, \muG)$.
\end{enumerate}
When $F(\sigma) \ge 0$, a constructive gauge achieving~(ii) is given by:
\[ G_\sigma^{\star}(p,q) := -\Ksym(p,q) + \frac{F(\sigma)}{P(\sigma)^2 \|\muG\|_2^4}\;\muG(p)\,\muG(q). \]
\end{theorem}
\begin{proof}
\textbf{(ii) $\Rightarrow$ (i):} If $\widetilde{K}_\sigma$ is PSD, evaluating the quadratic form on $\mathbf{1}$ gives $\langle \mathbf{1}, \widetilde{T}_\sigma \mathbf{1} \rangle_\sigma \ge 0$. By gauge invariance (the admissibility condition ensures the gauge contributes zero), this equals $F(\sigma)/P(\sigma)^2 \ge 0$.

\textbf{(i) $\Rightarrow$ (ii):} Assume $F(\sigma) \ge 0$. Define $G_\sigma^{\star}$ as above. Symmetry is inherited from $\Ksym$ and the rank-one term. The gauge condition holds: $\sum_{p,q} G_\sigma^{\star}\,\muG(p)\,\muG(q) = -F(\sigma)/P(\sigma)^2 + F(\sigma)/P(\sigma)^2 = 0$. The corrected kernel $\widetilde{K}_\sigma = c_\sigma\,\muG(p)\,\muG(q)$ with $c_\sigma = F(\sigma)/(P(\sigma)^2\|\muG\|_2^4) \ge 0$ is rank-one PSD.
\end{proof}

\begin{proposition}[S-Procedure Certificate]\label{prop:s-cert}
Let $f_1, \ldots, f_r$ be mean-zero functions under $\muG$. If there exist multipliers $\tau_1, \ldots, \tau_r \ge 0$ such that
\[ T_\sigma - \sum_{k=1}^{r} \tau_k\,|f_k\rangle\langle f_k| \;\ge\; 0 \qquad (\text{as operator on } \ell^2(\muG)), \]
then $F(\sigma) \ge 0$.
\end{proposition}
\begin{proof}
Each $|f_k\rangle\langle f_k|$ with $\EE_{\muG}[f_k] = 0$ defines an admissible gauge. Setting $G_\sigma = -\sum_k \tau_k |f_k\rangle\langle f_k|$, the corrected kernel $\Ksym + G_\sigma$ is PSD by hypothesis, and \Cref{thm:gauge-equiv} yields $F(\sigma) \ge 0$.
\end{proof}

\begin{remark}[Scope of the Certificate]
\Cref{thm:gauge-equiv} is a tautological equivalence at each fixed $\sigma$: $F(\sigma) \ge 0$ if and only if a PSD gauge exists. Since $F(\sigma) < 0$ for $\sigma > 1.14$ (see Part II), the gauge certificate is applicable only in the boundary layer $\sigma \in (1, \sigma_0)$. The global Riemann Hypothesis \emph{cannot} be reduced to $F(\sigma) \ge 0$ for all $\sigma > 1$, because this condition fails independently of RH.
\end{remark}

\section{Conclusion}

We have analyzed the spectral structure of the symmetrized arithmetic kernel and developed the S-procedure certificate framework. The stable negative inertia index of $2$ is a robust numerical feature, and the gauge equivalence theorem provides a clean tautological reformulation for each fixed $\sigma$. However, the sign change of $F(\sigma)$ at $\sigma_0 \approx 1.14$ prevents a direct global reduction of the Riemann Hypothesis to gauge positivity. The gap between the local prime-sum analysis and the global Li positivity remains the central open problem of the FST-RH program. The correct framework for addressing it is the Bombieri--Lagarias decomposition via $\log\xi$ (Part II, Theorem~4.1), not the Gibbs-normalized prime-sum formulation.

\appendix
\section{Spectral Refinement}

Numerical evaluation with $N=5000$ primes at $\sigma=1.01$ yields the following spectral structure:
\begin{enumerate}
    \item \textbf{Overlap:} The dominant negative eigenvector $v_1$ has a projection norm of $0.9431$ onto the feature space $W_\sigma = \operatorname{span}\{\log p, \log^2 p\}$.
    \item \textbf{Residuals:} The "tail" of the operator after removing the 2D feature projection has a minimum eigenvalue of $-4.8 \times 10^{-3}$, which is nearly an order of magnitude smaller than the raw $\lambda_1$.
    \item \textbf{Stability:} The index and overlap remain stable for $\sigma \in [1, 1.5]$, confirming that the instability is non-chaotic and concentrated in the logarithmic moments of the prime distribution.
\end{enumerate}

These results confirm that the global Gauge Positivity problem is essentially a stability requirement on the 2D coarse-grained arithmetic features, with the high-frequency tail being inherently more stable.

\bibliographystyle{plain}
\bibliography{fst-rh-references}

\end{document}
